% Lattice_Perturbative_Kappa.tex
% Investigation: Can kappa be derived from beta via perturbative lattice methods?
% Date: 2026-03-22
% Status: DEFINITIVE NEGATIVE RESULT -- perturbative route fails; confirms non-perturbative origin

\documentclass[11pt]{article}
\usepackage{amsmath,amssymb,amsfonts}
\usepackage{geometry}
\geometry{margin=1in}
\usepackage{booktabs}
\usepackage{hyperref}

\title{Perturbative Lattice Route to $\kappa$: \\
Why It Fails and What This Tells Us}
\author{DFD Research Note}
\date{March 22, 2026}

\begin{document}
\maketitle

%% ================================================================
\section{The Question}
%% ================================================================

The Ab Initio paper~\cite{Alcock2025} identifies $\beta_{U(1)} = 3.80$ as the
bare lattice coupling (Eq.~16) and measures the renormalized stiffness
$\kappa_{U(1)} \approx 7.25$ via the background-field method (Eq.~20).
The relationship
\begin{equation}\label{eq:ratio}
  \frac{\kappa}{\beta} \;\approx\; 1.91
\end{equation}
is stated to be ``non-perturbative and lattice-size dependent'' (p.~7,
Eq.~11).  This note investigates whether \emph{any} perturbative
expansion---naive, tadpole-improved, or boosted---can reproduce
Eq.~\eqref{eq:ratio}.

\medskip\noindent
\textbf{Verdict:} No.  Every perturbative method predicts $\kappa < \beta$,
whereas DFD measures $\kappa \approx 2\beta$.  The enhancement is
genuinely non-perturbative and arises from the microsector coupling
in the DFD action.


%% ================================================================
\section{DFD Setup Extracted from the Ab Initio Paper}
%% ================================================================

The DFD lattice action (Eq.~19 of \cite{Alcock2025}) is:
\begin{equation}\label{eq:DFD_action}
  S = \sum_x \bigl[-\log w(k_x)\bigr]
    + \frac{K_\psi}{2}\sum_{\langle xy\rangle}(\psi_x - \psi_y)^2
    - \sum_p \beta\, e^{-\psi_p}\cos(\theta_p + \theta_{\mathrm{bg}}\,\Omega_p),
\end{equation}
where:
\begin{itemize}
  \item $w(k) = \frac{2}{k+2}\sin^2\!\bigl(\frac{\pi}{k+2}\bigr)$ is the
        $SU(2)_k$ Chern--Simons microsector weight (Eq.~5).
  \item $\psi_x = \psi(k_x)$ is the coarse-graining map from the integer
        microsector field $k_x$ to the scalar density field.
  \item $\theta_p$ is the compact $U(1)$ plaquette angle (Wilson action).
  \item $\theta_{\mathrm{bg}}$ is the background twist; $\Omega_p$ marks
        the background-field plaquettes.
  \item $K_\psi = 0.25$, $k_0 = 8$ are auxiliary simulation parameters
        (verified to be irrelevant).
\end{itemize}

\noindent\textbf{Key point:} This is \emph{not} a standard compact $U(1)$
lattice gauge theory.  The gauge plaquette is modulated by $e^{-\psi_p}$,
coupling it to the microsector field.  The stiffness $\kappa$ therefore
receives contributions from \emph{both} the gauge sector and the
microsector fluctuations.

The stiffness is extracted via (Eq.~20):
\begin{equation}\label{eq:kappa_def}
  \kappa = \frac{F''(0)}{V}, \qquad
  F''(0) = \langle S''\rangle - \langle (S')^2\rangle + \langle S'\rangle^2,
\end{equation}
where primes denote derivatives with respect to $\theta_{\mathrm{bg}}$
at $\theta_{\mathrm{bg}} = 0$, and $V = L^3$.

The lattice parameters are:
\begin{center}
\begin{tabular}{ll}
\toprule
$\beta_{U(1)}$ & $3.80$ (input) \\
$\beta_{SU(2)}$ & $22.80$ (input, ratio $= 6$) \\
$\kappa_{U(1)}$ & $\approx 7.25$ (measured output) \\
$\kappa_{SU(2)}$ & $\approx 14.5$ (measured output) \\
$\kappa_{U(1)}/\kappa_{SU(2)}$ & $0.495 \pm 0.020$ (measured, predicted $= 1/2$) \\
$\alpha_W$ & $0.007297$ (measured, physical $= 0.0072974$) \\
\bottomrule
\end{tabular}
\end{center}

\noindent\textbf{Wilson vs Villain:} DFD uses the \emph{Wilson action}
$-\beta\cos\theta_p$, not the Villain action
$-\frac{\beta}{2}(\theta_p \bmod 2\pi)^2$.
For $\beta = 3.80 \gg \beta_c \approx 1.01$, both actions are
deep in the Coulomb phase and give equivalent universal physics,
differing only by $O(1/\beta^2)$ lattice artifacts.
The Villain model is Gaussian (exactly solvable on a torus after
summing over winding numbers), but does not apply directly to the
DFD microsector-coupled action.


%% ================================================================
\section{Standard Perturbative Expansion for Compact $U(1)$}
%% ================================================================

\subsection{Gaussian (One-Loop) Approximation}

For \emph{standard} compact $U(1)$ with Wilson action
$S = -\beta\sum_p \cos\theta_p$ on a 4D hypercubic lattice,
expanding $\cos\theta_p \approx 1 - \theta_p^2/2$ and performing
the Gaussian integral with the free lattice photon propagator
in Feynman gauge,
\[
  D_{\mu\nu}(k) = \frac{\delta_{\mu\nu}}{\beta\,\hat{k}^2}, \qquad
  \hat{k}^2 = \sum_\mu 4\sin^2(k_\mu/2),
\]
yields the one-loop plaquette expectation value:
\begin{equation}\label{eq:P_oneloop}
  \langle P\rangle = \langle\cos\theta_p\rangle
  = 1 - \frac{1}{2\beta} + O(1/\beta^2).
\end{equation}
At $\beta = 3.80$: $\langle P\rangle_{\text{1-loop}} = 0.8684$.

The coefficient $1/2$ arises because the plaquette field-strength
variance is $\langle\theta_p^2\rangle = 1/\beta$ (from the momentum
integral over the free propagator, using the identity
$\sum_\mu \hat{k}_\mu^2/\hat{k}^2 = 1$ by hypercubic symmetry).

\subsection{Higher-Order Perturbative Coefficients}

The weak-coupling expansion of the plaquette for compact $U(1)$ in 4D:
\begin{equation}\label{eq:P_series}
  \langle P\rangle = 1 - \frac{c_1}{\beta} - \frac{c_2}{\beta^2}
  - \frac{c_3}{\beta^3} - \cdots
\end{equation}
with known/estimated coefficients:
\begin{center}
\begin{tabular}{cll}
\toprule
$n$ & $c_n$ & Source \\
\midrule
1 & $1/2 = 0.500$ & Exact (free propagator integral) \\
2 & $\sim 0.125$ & Single-link Bessel expansion; vertex corrections \\
3 & $\sim 0.125$ & Estimated from $I_1(\beta)/I_0(\beta)$ asymptotics \\
\bottomrule
\end{tabular}
\end{center}
The single compact $U(1)$ link integral $\langle\cos\theta\rangle
= I_1(\beta)/I_0(\beta)$ has the large-$\beta$ expansion
$1 - \frac{1}{2\beta} - \frac{1}{8\beta^2} - \frac{1}{8\beta^3}
- \frac{25}{128\beta^4} - \cdots$\,.
At $\beta = 3.80$ the exact single-link value is
$I_1(3.80)/I_0(3.80) = 0.8554$.

The plaquette on a correlated lattice differs from the single-link result,
but the leading coefficient $c_1 = 1/2$ is exact and the higher coefficients
are of similar magnitude.


%% ================================================================
\section{Perturbative Stiffness: All Roads Lead to $\kappa < \beta$}
%% ================================================================

\subsection{Method 1: Direct Gaussian Calculation}

The helicity modulus (stiffness) is defined as
$\kappa = \partial^2 F/\partial\phi^2\big|_{\phi=0}$
where $\phi$ is a uniform boundary twist.  For the standard Wilson action:
\begin{align}
  F''(0)/V &= \beta\,\langle\cos\theta_p\rangle_{\Omega}
  - \beta^2\,\chi_{\sin}\;, \label{eq:helmod}
\end{align}
where $\langle\cdots\rangle_\Omega$ is the plaquette average over the
twisted direction and $\chi_{\sin}$ is the susceptibility of $\sin\theta_p$.

In the Gaussian approximation on an infinite lattice:
\begin{itemize}
  \item $\langle\cos\theta_p\rangle \approx 1 - 1/(2\beta)$
  \item $\chi_{\sin} \to 0$ as $V \to \infty$ (fluctuation term vanishes)
\end{itemize}
Therefore:
\begin{equation}\label{eq:kappa_gauss}
  \boxed{\kappa_{\text{Gaussian}} = \beta - \frac{1}{2} = 3.30}
  \qquad\text{at }\beta = 3.80.
\end{equation}
This is \emph{smaller} than $\beta$, giving $\kappa/\beta = 0.868$.

\subsection{Method 2: Tadpole Improvement (Lepage--Mackenzie)}

The Lepage--Mackenzie prescription~\cite{LepageMackenzie1993} identifies
the ``tadpole factor'' $u_0 = \langle P\rangle^{1/4}$ and defines the
improved coupling:
\begin{equation}
  \beta_I = \beta\,u_0^4 = \beta\,\langle P\rangle.
\end{equation}
The physical interpretation: each link variable $U_\ell$ receives a
mean-field renormalization $U_\ell \to u_0\,e^{ig\widetilde{A}_\mu}$,
absorbing the leading tadpole contribution.

For the stiffness, the tadpole-improved prediction is:
\begin{equation}\label{eq:kappa_TI}
  \kappa_{\text{TI}} \approx \beta\,\langle P\rangle \approx 3.42
  \qquad\text{(using $\langle P\rangle \approx 0.90$)}.
\end{equation}
This is \emph{still smaller} than $\beta$, with $\kappa_{\text{TI}}/\beta \approx 0.90$.

\medskip
\noindent\textbf{The fundamental issue:} The Lepage--Mackenzie improvement
was designed to bring perturbative predictions \emph{closer to MC data}
for quantities like Wilson loops and heavy-quark potentials.  It replaces
$g_0^2 \to g_0^2/u_0^4$, which \emph{increases} the effective coupling
constant (decreases the effective $\beta$).  For stiffness, this means
$\kappa_{\text{improved}} < \kappa_{\text{bare}} < \beta$.  Tadpole
improvement moves $\kappa$ in the \emph{wrong direction} relative to
what DFD measures.

\subsection{Method 3: Boosted Coupling}

One might try the ``boosted coupling'' $\beta_B = \beta/\langle P\rangle$,
which is larger than $\beta$.  This gives:
\begin{equation}
  \beta_B = \frac{3.80}{0.90} \approx 4.22.
\end{equation}
While this is larger than $\beta$, it is still far short of
$\kappa \approx 7.25$.  Moreover, the boosted coupling is the
\emph{inverse} of the Lepage--Mackenzie prescription and corresponds
to using $g^2 u_0^4$ rather than $g^2/u_0^4$ --- it is not a
physically motivated improvement scheme.

\subsection{Summary Table: All Perturbative Predictions}

\begin{center}
\begin{tabular}{lcc}
\toprule
Method & $\kappa$ prediction & $\kappa/\beta$ \\
\midrule
Tree level & $3.80$ & $1.000$ \\
Gaussian (1-loop) & $3.30$ & $0.868$ \\
2-loop estimate & $3.27$ & $0.860$ \\
Tadpole-improved & $3.42$ & $0.900$ \\
Boosted coupling & $4.22$ & $1.111$ \\
\midrule
\textbf{DFD measurement} & $\mathbf{7.25}$ & $\mathbf{1.908}$ \\
\bottomrule
\end{tabular}
\end{center}

\noindent Every perturbative method predicts $\kappa/\beta \leq 1.11$.
DFD measures $\kappa/\beta \approx 1.91$, a factor of $\sim\!2.2$ above
the Gaussian prediction.


%% ================================================================
\section{Why the Enhancement Is Non-Perturbative}
%% ================================================================

The factor $\kappa/\beta \approx 1.91$ cannot arise from any expansion
in $1/\beta$ applied to \emph{standard} compact $U(1)$.  Three distinct
non-perturbative mechanisms contribute:

\subsection{Microsector Coupling}

The DFD action~\eqref{eq:DFD_action} modulates the gauge plaquette
by $e^{-\psi_p}$, where $\psi_p$ depends on the microsector field $k_x$.
The stiffness~\eqref{eq:kappa_def} therefore receives contributions from
microsector fluctuations.  Specifically, the second derivative of the
action with respect to $\theta_{\mathrm{bg}}$ includes terms of the form
\[
  \beta^2 \bigl\langle (e^{-\psi_p}\sin\theta_p)^2 \bigr\rangle,
\]
which involve correlated gauge-microsector fluctuations absent in
standard $U(1)$ theory.  If $e^{-\psi_p}$ fluctuates, Jensen's inequality
gives $\langle e^{-\psi}\rangle > e^{-\langle\psi\rangle}$, providing
a mechanism for \emph{enhancing} the effective coupling.

\subsection{Topology of $\mathbb{C}P^2$}

The DFD internal manifold is $\mathbb{C}P^2 \times S^3$, not a 4-torus.
The non-trivial topology of $\mathbb{C}P^2$ (Euler characteristic
$\chi = 3$, non-trivial second cohomology) affects the gauge field
dynamics through:
\begin{itemize}
  \item Monopole structure: compact $U(1)$ on $\mathbb{C}P^2$ has different
        topological sectors than on a torus.
  \item Mode counting: the number of zero modes and near-zero modes
        of the lattice Laplacian differs.
  \item Curvature effects: the curvature of $\mathbb{C}P^2$ modifies
        the effective action.
\end{itemize}

\subsection{Finite-Size Enhancement}

The paper reports $\kappa$ values that depend on lattice size $L$
(Table~VI), with the physical $\alpha$ emerging across $L = 6, 8, 10, 12$
within $\sim\!1\%$.  Finite-volume effects on a compact manifold
can enhance the stiffness relative to the infinite-volume perturbative
prediction.


%% ================================================================
\section{Exact and Semi-Exact Results}
%% ================================================================

\subsection{Villain Model}

The Villain action
$S_V = -\frac{\beta}{2}\sum_p \min_n (\theta_p + 2\pi n)^2$
is Gaussian after summing over winding numbers and is exactly
solvable on a torus via Poisson resummation.  For the Villain model:
\begin{equation}
  \kappa_{\text{Villain}} = \beta \qquad\text{(exactly, on a torus)}.
\end{equation}
The Villain and Wilson actions agree at large $\beta$ up to $O(1/\beta^2)$
corrections from the $\cos\theta$ vs.\ $\theta^2/2$ difference.
At $\beta = 3.80$, these corrections are $\sim\!2\%$ and cannot produce
the observed factor-of-2 enhancement.

The Villain model on $\mathbb{C}P^2$ (as opposed to a torus) has not been
studied in the literature and would require separate analysis to determine
whether the non-trivial topology produces a similar enhancement.

\subsection{Single-Link Exact Result}

For a single compact $U(1)$ link with Boltzmann weight $e^{\beta\cos\theta}$:
\begin{equation}
  \langle\cos\theta\rangle = \frac{I_1(\beta)}{I_0(\beta)}
  = 0.8554 \quad\text{at }\beta = 3.80.
\end{equation}
This provides the mean-field tadpole factor $u_0 = 0.8554^{1/4} = 0.9617$.


%% ================================================================
\section{The Required Value of $\kappa$ and Its Origin}
%% ================================================================

From the electroweak mixing relation (Eq.~13 of \cite{Alcock2025}):
\begin{equation}
  \alpha_W = \frac{(1/\kappa_{U(1)})(4/\kappa_{SU(2)})}
             {(1/\kappa_{U(1)}) + (4/\kappa_{SU(2)})} \cdot \frac{1}{4\pi}.
\end{equation}
Using the Frame Stiffness Theorem $\kappa_{U(1)}/\kappa_{SU(2)} = 1/2$:
\begin{equation}
  \alpha_W = \frac{1}{6\pi\,\kappa_{U(1)}}.
\end{equation}
Setting $\alpha_W = \alpha_{\text{phys}} = 0.0072974$:
\begin{equation}\label{eq:kappa_required}
  \boxed{\kappa_{U(1)} = \frac{1}{6\pi\alpha_{\text{phys}}} = 7.270}
\end{equation}
giving $\kappa/\beta = 7.270/3.80 = 1.913$.

\medskip
\noindent\textbf{This value is determined by the lattice simulation,
not by perturbation theory.}  The fact that the simulation produces
$\kappa \approx 7.27$ at $\beta = 3.80$ is the central numerical
prediction of the Ab Initio paper; the perturbative route is a dead end.


%% ================================================================
\section{Conclusions}
%% ================================================================

\begin{enumerate}
\item \textbf{Perturbative lattice methods fail:} The Gaussian approximation
  gives $\kappa = \beta - 1/2 = 3.30$.  Tadpole improvement gives
  $\kappa \approx 3.42$.  All perturbative methods predict $\kappa < \beta$,
  whereas DFD measures $\kappa \approx 2\beta$.

\item \textbf{The enhancement is a factor of $\sim\!2.2$} over the Gaussian
  prediction, confirming the paper's statement that the $\beta$-to-$\kappa$
  relationship is non-perturbative.

\item \textbf{The microsector coupling is the key:} The DFD action
  is not standard compact $U(1)$---it couples gauge fields to the
  microsector via $e^{-\psi_p}$.  This coupling generates the
  stiffness enhancement that produces $\alpha = 1/137$.

\item \textbf{No analytic shortcut exists:} The value $\kappa/\beta \approx 1.91$
  can only be obtained by Monte Carlo simulation of the full DFD action.
  This is precisely the role of the lattice calculation in the derivation chain.

\item \textbf{Wilson vs Villain:} DFD uses the Wilson action.
  The Villain action is exactly solvable on a torus (giving $\kappa = \beta$
  exactly), but neither action on a torus can produce $\kappa/\beta \approx 1.91$.
  The enhancement requires the full DFD microsector-coupled action.

\item \textbf{The spectral action $a_4$ route} (previously investigated)
  attempts to derive $\kappa$ analytically from the Seeley--DeWitt
  coefficients on $\mathbb{C}P^2 \times S^3$.  The perturbative lattice
  route investigated here is complementary but equally unable to bypass
  the Monte Carlo step.
\end{enumerate}


\begin{thebibliography}{9}

\bibitem{Alcock2025}
G.~Alcock,
``Ab Initio Derivation of the Fine Structure Constant from Density Field Dynamics,''
(December 27, 2025).

\bibitem{LepageMackenzie1993}
G.~P.~Lepage and P.~B.~Mackenzie,
``On the Viability of Lattice Perturbation Theory,''
Phys.\ Rev.\ D \textbf{48}, 2250 (1993)
[arXiv:hep-lat/9209022].

\bibitem{Creutz1983}
M.~Creutz,
\emph{Quarks, Gluons and Lattices}
(Cambridge University Press, 1983).

\bibitem{Wilson1974}
K.~G.~Wilson,
``Confinement of Quarks,''
Phys.\ Rev.\ D \textbf{10}, 2445 (1974).

\bibitem{HelicityModulus}
L.~Polley and U.-J.~Wiese (Eds.),
``The helicity modulus in gauge field theories,''
arXiv:hep-lat/0311007.

\bibitem{VillainAction}
J.~Villain,
``Theory of one- and two-dimensional magnets with an easy
magnetization plane,''
J.\ Physique \textbf{36}, 581 (1975).

\bibitem{ImprovedSpinWave}
S.~Chatterjee,
``Improved Spin-Wave Estimate for Wilson Loops in $U(1)$ Lattice Gauge Theory,''
arXiv:2106.00705.

\end{thebibliography}

\end{document}
